\section{Architecture}
\label{sec:architecture}

\subsection{The three layers}

The Week 8 lecture's layers map directly: the React application is the presentation layer, the
FastAPI service is the application layer, and the data layer is read-only files plus one small
SQLite file. Figure~\ref{fig:arch} is also the first diagram of the report.

\begin{figure}[H]
\centering
\resizebox{\textwidth}{!}{%
\begin{tikzpicture}[font=\small,
  box/.style={draw=ink, rounded corners=2pt, align=center, minimum width=3.7cm, minimum height=9mm, inner sep=3pt},
  fe/.style={box, fill=accentsoft},
  be/.style={box, fill=tealsoft},
  ml/.style={box, fill=violetsoft},
  st/.style={box, fill=greysoft},
  arr/.style={-{Stealth[length=2.2mm]}, thick, draw=ink},
  darr/.style={{Stealth[length=2.2mm]}-{Stealth[length=2.2mm]}, thick, draw=ink}]
  % front end
  \node[fe] (views)  at (0, 0)    {Views V1--V8\\ \footnotesize screens, forms, feedback states};
  \node[fe] (charts) at (0,-1.5)  {Charts G1--G9\\ \footnotesize one Plotly.js wrapper};
  \node[fe] (valid)  at (0,-3.0)  {Validation\\ \footnotesize address checksum, what-if rules};
  \node[fe] (client) at (0,-4.5)  {API client\\ \footnotesize Axios, \texttt{ApiError}, request cache};
  % back end
  \node[be] (routers)  at (6.2, 0)   {Routers\\ \footnotesize \texttt{/api}: E1--E22};
  \node[be] (schemas)  at (6.2,-1.5) {Schemas\\ \footnotesize Pydantic: requests, responses, errors};
  \node[be] (services) at (6.2,-3.0) {Services\\ \footnotesize explorer, what-if, export, assistant, live};
  \node[be] (core)     at (6.2,-4.5) {Core\\ \footnotesize settings, error handlers, logging};
  % model layer and stores
  \node[ml] (explorer) at (12.4, 0)   {\texttt{explorer} package (A2)\\ \footnotesize \texttt{Predictor}, features, methods, explain};
  \node[st] (files)    at (12.4,-1.5) {Models and snapshot tables\\ \footnotesize \texttt{.joblib}, \texttt{.csv}, \texttt{settings.json}};
  \node[st] (db)       at (12.4,-3.0) {SQLite file\\ \footnotesize watchlist, live cache index};
  \node[st, dashed] (opt) at (12.4,-4.5) {Optional\\ \footnotesize Assistant weights; Esplora API};
  % groups
  \begin{scope}[on background layer]
    \node[draw=accent, rounded corners=4pt, inner sep=9pt, fit=(views)(client), label={[font=\small\bfseries, text=accent]above:React front end (browser)}] (fegrp) {};
    \node[draw=teal, rounded corners=4pt, inner sep=9pt, fit=(routers)(core), label={[font=\small\bfseries, text=teal]above:FastAPI back end (Uvicorn)}] (begrp) {};
    \node[draw=violet, rounded corners=4pt, inner sep=9pt, fit=(explorer)(opt), label={[font=\small\bfseries, text=violet]above:Model and data}] (mlgrp) {};
  \end{scope}
  % arrows
  \draw[arr] (views) -- (charts);
  \draw[arr] (valid) -- (client);
  \draw[darr] (client.east) -- node[above, font=\footnotesize, align=center] {HTTP and JSON} node[below, font=\footnotesize, align=center] {GET POST PUT DELETE} (core.west |- client.east);
  \draw[arr] (routers) -- (schemas);
  \draw[arr] (schemas) -- (services);
  \draw[arr] (services.east) -- (explorer.west);
  \draw[arr] (explorer) -- (files);
  \draw[arr] (services.east) -- (db.west);
  \draw[arr, dashed] (services.east) -- (opt.west);
\end{tikzpicture}}
\caption{The layers. Solid arrows are calls made in the Must tier. Dashed parts are optional: the
application runs without them and says what is missing.}
\label{fig:arch}
\end{figure}

\textbf{Rules of the layering.}
\begin{itemize}
  \item A router only parses the request, calls one service function and returns its result. It
    never touches pandas or a model.
  \item A service never knows about HTTP. It raises the package's own exceptions
    (\code{Unsupported}, \code{InsufficientHistory}, \code{ValueError}); the core layer turns them
    into responses in one place.
  \item The \code{explorer} package is Assignment 2's code, copied unchanged. Every difference is
    in a service around it and is listed in Section~\ref{sec:mods}.
  \item A view never calls Axios. It calls a typed function of the API client and renders one of
    the request's four states: idle, loading, success, error.
  \item A view never draws with Plotly directly. It passes data to a chart component, which uses
    the one wrapper.
\end{itemize}

\subsection{One request, end to end}
\label{sec:walkthrough}

The what-if is the path the report and the video explain, because it crosses every layer.

\begin{enumerate}
  \item The user changes the amount in view V4. The form's rules run on the keystroke; an invalid
    value shows its message under the field and nothing is sent.
  \item 250 ms after the last change, the client cancels any request still in flight and sends
    \code{POST /api/addresses/{address}/what-if} with \code{amountBtc}, \code{direction} and
    \code{time}.
  \item FastAPI parses the body into a Pydantic model. A wrong type, an unknown field or a value
    out of range answers 400 with the \code{invalid} body and the failing field.
  \item The what-if service checks the address (404), converts the amount to whole satoshis and
    calls \code{Predictor.score_what_if}.
  \item \emph{Preprocessing.} The package builds the hypothetical row, giving it the median shape of
    the earlier transactions, appends it to the history strictly before its time, and computes the
    features F1 to F4 and their robust z-scores against the address's own comparison windows.
  \item \emph{Prediction.} The clipped z-scores and the direction go to each saved model; each
    returns a score, compared with that model's own threshold.
  \item \emph{Postprocessing.} The result becomes a transaction detail: values in BTC, the role of
    each feature, the comparison basis, every model's score and flag, the template explanation and
    the limitation line.
  \item The response model serialises it. The client resolves; the verdict card, the feature chart
    and the model table re-render, and the curve request (E9) follows.
\end{enumerate}

\subsection{Technology choices}
\label{sec:stack}

Versions are pinned on the first day and recorded in \code{environment.yml} and
\code{package-lock.json}; the report quotes them from those files.

{\small
\begin{longtable}{@{}L{2.5cm}L{3.3cm}L{8.2cm}@{}}
\toprule
\textbf{Part} & \textbf{Choice} & \textbf{Why, and what was not chosen} \\
\midrule
\endhead
Interface & React with TypeScript & Required by the brief; TypeScript carries the API contract into every screen (D-A3-2) \\
Build tool & Vite & Starts in seconds, proxies \code{/api}; create-react-app is deprecated \\
Charts & Plotly.js & D-A3-3. Not D3.js (every interaction by hand) or Chart.js (not on the High Distinction line) \\
Requests & Axios & D-A3-4. Cancellation and time limits for the live what-if \\
Routing & The prototype's hash router & Every view and state has a link; no server rewrite rules. Not React Router: 60 lines already do the job and are tested \\
Server state & The prototype's \code{useRequest} hook, extended with cancellation and polling & One small hook the team can explain. Not a query library: one more dependency for what three additions provide \\
Styling & The prototype's CSS with design tokens & Already responsive, themed and contrast-checked. Not a component kit: restyling every screen would gain nothing the rubric marks \\
Server & FastAPI on Uvicorn & Required by the brief; typed routes; generated OpenAPI documentation \\
Validation & Pydantic models & One definition gives parsing, validation, the response shape and the documentation \\
Model layer & Assignment 2's \code{explorer} package, pandas, scikit-learn, joblib & The brief: integrate the Assignment 2 model. Versions equal to Assignment 2's (D-A3-10) \\
Storage & CSV and joblib files; SQLite through the standard library & D-A3-8. Not an ORM: three small tables \\
Back-end tests & pytest with FastAPI's test client & Assignment 2's tool; requests run in-process, no server needed \\
Front-end tests & Vitest and Testing Library & The prototype's 66 tests keep running \\
Browser tests & Playwright (S-11) & Three widths, real browsers, screenshots for the report \\
\bottomrule
\end{longtable}
}

\subsection{Repository layout}
\label{sec:layout}

\begin{lstlisting}
COS30049A3/
  README.md                    set-up, run, configuration, tests, troubleshooting
  LICENSES.md  CHANGELOG.md
  backend/
    environment.yml  requirements.txt  requirements-assistant.txt
    app/
      main.py                  create_app(): lifespan, middleware, handlers, routers
      api/                     one router per resource
        system.py              E1 health, E2 config
        addresses.py           E3 list, E4 lookup, E5 timeline, E6 transactions, E20 delete
        transactions.py        E7 detail
        whatif.py              E8 single, E9 curve
        exports.py             E10
        methods.py             E11
        explanations.py        E12
        watchlist.py           E13-E15
        live.py                E16-E19
      schemas/                 Pydantic models: common, addresses, transactions, whatif, errors ...
      services/
        store.py               AddressStore: snapshot (and live) data behind one interface
        explorer_service.py    timeline, rows, detail; the per-address cache; method override
        whatif_service.py      single and batched what-if
        export_service.py      CSV and JSON
        assistant_service.py   lazy load, one answer at a time, cache
        watchlist_repo.py      SQLite
        live_service.py        jobs, Esplora fetch, progress events
      core/
        config.py              environment variables and settings.json -> web
        errors.py              exception classes and the handlers that build the error body
        validation.py          Bitcoin address and transaction-ID checks
        logging.py
    explorer/  lm/             Assignment 2's runtime modules, copied unchanged (A2_VENDOR.json)
    config/settings.json       Assignment 2's settings plus the "web" block
    models/                    the saved methods (about 30 MB)
    data/processed/esplora/    features.csv, transactions.csv, scores.csv, txgpt_scores.csv, manifest.json
    var/                       explorer.db, live cache, logs (git-ignored, created at start)
    tests/                     unit, API, parity, wording, performance smoke
    MODIFICATIONS.md           every change made for the web, with its test
  frontend/
    package.json  package-lock.json  vite.config.ts  tsconfig.json  index.html
    src/
      api/                     http.ts (Axios), endpoints.ts, types.ts, mock/ (tests only)
      views/                   V1-V8, one folder each
      charts/                  PlotlyChart.tsx (the wrapper) and G1-G9
      components/              buttons, badges, table, modal, toast, tabs, filter bar, fields ...
      hooks/                   useRequest, useDebounced, usePolling, useEventSource, usePreferences
      lib/                     address.ts, dates.ts, format.ts, router.ts, validation/
      styles/
    tests/                     unit, component, journeys
    e2e/                       Playwright (S-11)
  tools/
    vendor_a2.py               copies the runtime from the Assignment 2 repository, with hashes
    package.py                 builds the zip from a git tag, with SHA256SUMS.txt
    measure.py                 the performance table (S-12)
  docs/
    plan/  A3_delivery_plan.pdf   this plan
    brief/                     the brief (not in the zip)
    AI_USE_LOG.md  TEST_LOG.md  change-log.md  usability/
  report/                      the report's LaTeX (not in the zip)
\end{lstlisting}

\subsection{The address store}
\label{sec:store}

Services read addresses through one interface, so the snapshot and a fetched address look the same
to every endpoint.

\begin{itemize}
  \item \code{AddressStore.get(address)} returns the address's features, transactions and scores
    and its \code{source}: \code{snapshot} or \code{live}. It raises \code{Unsupported} when the
    address is in neither.
  \item The snapshot store wraps the \code{Predictor}'s three tables. It is filled once at start
    and never changes.
  \item The live store (S-09) adds an address after a lookup has fetched, cleaned, featured and
    scored its history. It is a folder of cached tables under \fp{var/live/}, named by the SHA-256
    of the address so that no user text reaches a file path.
  \item Every response that shows data carries \code{source} and its date, so the interface can
    label live data as live.
\end{itemize}

\subsection{Configuration}
\label{sec:config}

No value is typed into code. Numbers the web tier needs go into a new \code{web} block of
\fp{config/settings.json}; the front end reads them from E2, so the two tiers cannot disagree.

\begin{lstlisting}[style=json]
"web": {
  "whatIf":   { "minAmountBtc": 0.00000001, "maxAmountBtc": 21000000, "maxGapYears": 5,
                "curvePoints": 41, "curveMaxPoints": 61, "debounceMs": 250 },
  "lists":    { "pageSize": 50, "maxPageSize": 200 },
  "requests": { "timeoutMs": 10000, "assistantTimeoutMs": 90000, "healthPollSeconds": 15 },
  "charts":   { "webglMinPoints": 1000 },
  "watchlist":{ "maxEntries": 200, "maxLabelChars": 40, "maxNoteChars": 500 },
  "live":     { "defaultMaxTx": 500, "maxTxLimit": 1000, "lookupsPerHour": 6, "cacheDays": 7 },
  "replay":   { "speeds": [1, 5, 20] },
  "limits":   { "maxBodyBytes": 10240, "maxAddressChars": 100 }
}
\end{lstlisting}

Values here are starting points\est{}; a change is logged in \fp{CHANGELOG.md} as in Assignment 2.
Environment variables only choose what a machine may do:

\begin{center}
\small
\begin{tabularx}{\textwidth}{@{}L{4.6cm}L{4.2cm}Y@{}}
\toprule
\textbf{Variable} & \textbf{Default} & \textbf{Meaning} \\
\midrule
\code{EXPLORER_CORS_ORIGINS} & \code{http://localhost:5173}, \code{http://127.0.0.1:5173} & Origins allowed to call the API directly \\
\code{EXPLORER_LIVE_LOOKUP} & \code{off} & Turns S-09 on; needs CR-08 \\
\code{EXPLORER_ASSISTANT} & \code{auto} & \code{auto}: on when the weights are present; \code{off}: template only \\
\code{EXPLORER_LOG_LEVEL} & \code{INFO} & Log detail \\
\code{VITE_API_BASE} (front end) & empty & Empty: same origin through Vite's proxy; else the back end's address \\
\code{VITE_API_MODE} (front end) & \code{http} & \code{mock} runs the screens on the in-browser mock, for tests \\
\bottomrule
\end{tabularx}
\end{center}

\subsection{Ways to run it}
\label{sec:runmodes}

\begin{center}
\small
\begin{tabularx}{\textwidth}{@{}L{2.8cm}Y L{1.6cm}@{}}
\toprule
\textbf{Mode} & \textbf{How} & \textbf{Tier} \\
\midrule
Development and marking & Terminal 1: \code{uvicorn app.main:app --port 8000} in \fp{backend/}. Terminal 2: \code{npm run dev} in \fp{frontend/}, then open \code{http://localhost:5173}. The API's documentation is at \code{http://localhost:8000/docs} & \must \\
Tests & \code{python -m pytest -q} in \fp{backend/}; \code{npm test} in \fp{frontend/}. Neither needs the other tier running & \must \\
One command & \code{npm run build}, then Uvicorn serves \fp{frontend/dist} at \code{/}; or \code{docker compose up} & \could{} (C-05) \\
\bottomrule
\end{tabularx}
\end{center}

\begin{done}
A fresh clone starts both tiers with the README's commands in under 15 minutes including installs,
the back end reports ready within 10 seconds of starting\est{} (5.6 s measured for import and load
on the Mac), and the browser console and both terminals show no error or warning.
\end{done}
